\chapter{Подробные решения}

\section{З1--З4. Переводы}
\subsection*{З1. Из десятичной системы}\label{sol:z1}
\hyperref[task:z1]{К условию.}
Деления на 2 дают пары «частное, остаток»:
\[
(107,0),\ (53,1),\ (26,1),\ (13,0),\ (6,1),\ (3,0),\ (1,1),\ (0,1).
\]
Остатки снизу вверх: $11010110_2$.
Группировка по три бита: \texttt{011 010 110}, то есть $326_8$; по четыре: \texttt{1101 0110}, то есть $D6_{16}$.

Проверки: $128+64+16+4+2=214$;
$3\cdot64+2\cdot8+6=214$;
$13\cdot16+6=214$.

\subsection*{З2. Веса разрядов}\label{sol:z2}
\hyperref[task:z2]{К условию.}
\begin{align*}
1011011_2&=64+16+8+2+1=91,\\
273_8&=2\cdot64+7\cdot8+3=187,\\
7D_{16}&=7\cdot16+13=125.
\end{align*}
Цифра D имеет значение 13.

\subsection*{З3. Группы справа налево}\label{sol:z3}
\hyperref[task:z3]{К условию.}
\begin{itemize}
    \item \texttt{11 10 10 11 01} даёт $32231_4$.
    \item \texttt{001 110 101 101} даёт $1655_8$.
    \item \texttt{0011 1010 1101} даёт $3AD_{16}$.
\end{itemize}
Для проверки все записи дают 941: $3\cdot256+10\cdot16+13=941$.
Дополнялись только старшие группы; значение от этого не меняется.

\subsection*{З4. Сохранение внутренних нулей}\label{sol:z4}
\hyperref[task:z4]{К условию.}
Каждой hex-цифре сопоставляем ровно четыре бита:
\[
3C7_{16}\longrightarrow\texttt{0011 1100 0111}.
\]
Убираем только ведущие нули: $1111000111_2$. Для восьмеричной записи группируем заново: \texttt{001 111 000 111}, ответ $1707_8$.
Проверка: $3\cdot256+12\cdot16+7=967=512+7\cdot64+7$.

В $102_2$ запрещена цифра 2, в $19_8$ --- 9, в $2G_{16}$ --- G. В $203_4$ все цифры меньше 4, поэтому запись допустима (её значение равно 35).

\section{З5--З8. Арифметика}
\subsection*{З5. Переносы и заёмы}\label{sol:z5}
\hyperref[task:z5]{К условию.}
Сложение по разрядам справа: $1+0=1$; $1+1=10_2$; $0+1+1=10_2$; $1+0+1=10_2$; $1+1+1=11_2$. Получаем $110001_2$.

Для вычитания выравниваем запись: \texttt{101101} минус \texttt{001110}. Младший бит: $1-0=1$. В следующем разряде $0-1$ требует займа у бита 2; он становится нулём. Для бита 2 нужен новый заём у бита 3. Затем в бите 3 также нужен заём, который проходит через нулевой бит 4 от бита 5. В итоге остаётся \texttt{011111}, то есть $11111_2$.

Проверка: $27+22=49$; $45-14=31$.

\subsection*{З6. Основание 16}\label{sol:z6}
\hyperref[task:z6]{К условию.}
Сложение: $7+D=20_{10}=14_{16}$, пишем 4, перенос 1. Затем $B+8+1=20_{10}=14_{16}$, снова 4 и перенос 1. Наконец $2+0+1=3$. Ответ $344_{16}$.

Вычитание $300_{16}-0B6_{16}$: для младшего разряда занимаем через средний ноль у старшей цифры 3. Старшая становится 2, средняя после передачи единицы младшему --- F, младшая получает 16. Тогда $16-6=10=A$; $15-11=4$; $2-0=2$. Ответ $24A_{16}$.

Проверка: $695+141=836$; $768-182=586$.

\subsection*{З7. Частичные произведения}\label{sol:z7}
\hyperref[task:z7]{К условию.}
Для $34_8\cdot6_8$: $4\cdot6=24_{10}=30_8$ --- пишем 0, переносим 3. Далее $3\cdot6+3=21_{10}=25_8$ --- пишем 5, переносим 2. Ответ $250_8$. Проверка: $28\cdot6=168=2\cdot64+5\cdot8$.

Двоичное умножение:
\begin{verbatim}
      1011
   x   101
   -------
      1011
     0000
  + 1011
   -------
    110111
\end{verbatim}
Складываем ненулевые частичные произведения:
\[
1011_2+101100_2=110111_2.
\]
Проверка: $11\cdot5=55$.

\subsection*{З8. Деление с остатком}\label{sol:z8}
\hyperref[task:z8]{К условию.}
Для $325_8:6_8$ берём первые две цифры $32_8$. Произведение $6_8\cdot4_8=30_8$, остаток $2_8$. Сносим 5: получаем $25_8$. Следующая цифра частного 3, так как $6_8\cdot3_8=22_8$, остаток 3. Ответ: $43_8$, остаток $3_8$.

Для $2D9_{16}:B_{16}$ первые две цифры $2D_{16}$. Произведение $B_{16}\cdot4_{16}=2C_{16}$, остаток 1. Сносим 9: $19_{16}$. Следующая цифра 2, произведение $16_{16}$, остаток 3. Ответ: $42_{16}$, остаток $3_{16}$.

Десятичные проверки: $213=6\cdot35+3$, $3<6$;
$729=11\cdot66+3$, $3<11$.

\section{З9--З12. Диапазоны}
\subsection*{З9. Два нуля}\label{sol:z9}
\hyperref[task:z9]{К условию.}
$2^6=64$ --- количество комбинаций во всех четырёх форматах. В unsigned все комбинации дают разные значения $0\ldots63$.

Прямой и обратный коды имеют пять битов для диапазона модулей до $2^5-1=31$. Диапазон $[-31,31]$ содержит $31+1+31=63$ числа. У прямого кода нули \texttt{000000} и \texttt{100000}; у обратного --- \texttt{000000} и \texttt{111111}.

В дополнительном коде ноль один; диапазон $[-32,31]$ содержит $32+1+31=64$ значения.

\subsection*{З10. Границы в битах}\label{sol:z10}
\hyperref[task:z10]{К условию.}
Для unsigned максимум $2^9-1=511$, минимум 0: коды из девяти единиц и девяти нулей соответственно.

Для дополнительного кода минимум $-2^8=-256$ имеет код \texttt{100000000}; максимум $2^8-1=255$ --- \texttt{011111111}. В обоих форматах используется $2^9=512$ различных значений. Код из всех единиц в дополнительном коде означает $-1$.

\subsection*{З11. Доказательство минимальности}\label{sol:z11}
\hyperref[task:z11]{К условию.}
\begin{enumerate}
    \item 9 битов unsigned дают максимум 511, 10 битов --- 1023. Поскольку $511<1000\le1023$, нужны 10 битов.
    \item В 9 битах дополнительного кода $[-256,255]$ не покрывает $[-300,300]$. В 10 битах $[-512,511]$ покрывает. Ответ 10.
    \item В 8 битах $[-128,127]$ не помещается верхняя граница 128. В 9 битах $[-256,255]$ подходят обе. Ответ 9.
\end{enumerate}

\subsection*{З12. Проверка до кодирования}\label{sol:z12}
\hyperref[task:z12]{К условию.}
Сначала проверяем диапазоны, затем строим коды:
\begin{itemize}
    \item Unsigned $[0,255]$: $127=7F_{16}$, $128=80_{16}$; отрицательное число не помещается.
    \item Прямой и обратный $[-127,127]$: из трёх чисел помещается только 127 с кодом $7F_{16}$.
    \item Дополнительный $[-128,127]$: $-128$ имеет код $80_{16}$, 127 --- $7F_{16}$; $+128$ не помещается.
    \item Со смещением 128 диапазон также $[-128,127]$. Для $-128$ храним $-128+128=0$, код $00_{16}$; для 127 храним 255, код $FF_{16}$. Число 128 потребовало бы код 256, то есть девять битов.
\end{itemize}
Одинаковый диапазон дополнительного и смещённого кодов не означает одинаковые записи чисел.

\section{З13--З16. Представления}
\subsection*{З13. Положительное и отрицательное число}\label{sol:z13}
\hyperref[task:z13]{К условию.}
$46=32+8+4+2$, поэтому положительная запись \texttt{00101110}.
\begin{center}
\small
\begin{tabular}{l|c|c}
Код & $+46$ & $-46$\\\hline
Прямой & \texttt{00101110} (2E) & \texttt{10101110} (AE)\\
Обратный & \texttt{00101110} (2E) & \texttt{11010001} (D1)\\
Дополнительный & \texttt{00101110} (2E) & \texttt{11010010} (D2)\\
Со смещением 128 & \texttt{10101110} (AE) & \texttt{01010010} (52)
\end{tabular}
\end{center}
В прямом меняется знак; в обратном инвертируются все биты; в дополнительном к инверсии добавляется 1. Для смещённого вычисляем $46+128=174$ и $-46+128=82$. Обратная проверка последнего: $82-128=-46$.

\subsection*{З14. Интерпретация слова}\label{sol:z14}
\hyperref[task:z14]{К условию.}
$D6_{16}=\texttt{11010110}$; беззнаковое значение $13\cdot16+6=214$.
\begin{itemize}
    \item Прямой: знак 1, модуль \texttt{1010110}, то есть $64+16+4+2=86$. Значение $-86$.
    \item Обратный: инверсия \texttt{00101001} даёт модуль 41. Значение $-41$.
    \item Дополнительный: $214-256=-42$.
    \item Со смещением: $214-128=86$.
\end{itemize}
Нельзя определять значение только по старшему биту, не зная формата.

\subsection*{З15. Коррекция смещения}\label{sol:z15}
\hyperref[task:z15]{К условию.}
Коды: $u_A=-40+128=88=58_{16}$, $u_B=17+128=145=91_{16}$. Точная сумма значений равна $-23$ и помещается в диапазон $[-128,127]$.

Простая сумма кодов равна $88+145=233=E9_{16}$. В \emph{исходном смещённом формате} получаем $233-128=105$. Ожидаемая сумма равна $-23$, поэтому требуется коррекция смещения.

Коррекция: $u_{A+B}=88+145-128=105=69_{16}$.
Проверка: $105-128=-23$. Дополнительный код суммы: $256-23=233=E9_{16}$. В данном случае сырая сумма кодов совпала с дополнительным кодом ответа, но формат результата при таком чтении уже другой.

\subsection*{З16. Минимальное число}\label{sol:z16}
\hyperref[task:z16]{К условию.}
$-1$ кодируется $256-1=255$, то есть \texttt{11111111}; 0 --- \texttt{00000000}; $-128$ --- $256-128=128$, то есть \texttt{10000000}.

Инверсия \texttt{10000000} даёт \texttt{01111111}, прибавление 1 --- снова \texttt{10000000}. По правилу чтения это по-прежнему $-128$. Точный результат смены знака должен быть $+128$, но он превосходит максимум 127. При смене знака произошло знаковое переполнение.

\section{З17--З20. Арифметика в 8 битах}
\subsection*{З17. Сложение разных знаков}\label{sol:z17}
\hyperref[task:z17]{К условию.}
Коды $52=34_{16}$ и $-19=256-19=237=ED_{16}$.
Сумма кодов $52+237=289=121_{16}$. Сохраняем $21_{16}=\texttt{00100001}$, то есть 33 и в unsigned, и в дополнительном коде.

CF=1, поскольку $289\ge256$. Точная знаковая сумма $52-19=33$ представима, поэтому OF=0. Результат ненулевой с нулевым старшим битом: ZF=0, SF=0.

\subsection*{З18. Два знаковых переполнения}\label{sol:z18}
\hyperref[task:z18]{К условию.}
$95+48=143=8F_{16}=\texttt{10001111}$. Беззнаково 143, знаково $143-256=-113$. Переноса за байт нет: CF=0. Точная сумма 143 больше 127: OF=1. ZF=0, SF=1.

Для отрицательных слагаемых коды $9C_{16}$ (156) и $D8_{16}$ (216). Их сумма $372=174_{16}$, после усечения $74_{16}=\texttt{01110100}$, то есть 116 в обеих интерпретациях. CF=1. Точная сумма $-140<-128$: OF=1. ZF=0, SF=0.

Во втором случае сохранённый результат положителен, хотя точная сумма отрицательна: при OF=1 знак сохранённого слова не является знаком точного ответа.

\subsection*{З19. CF и OF}\label{sol:z19}
\hyperref[task:z19]{К условию.}
$F0_{16}+10_{16}=100_{16}$: сохранённый код $00_{16}$, обе интерпретации равны 0. CF=1. Знаково операнды равны $-16$ и 16, сумма 0 представима: OF=0. ZF=1, SF=0.

$7F_{16}+01_{16}=80_{16}$: unsigned 128, signed $-128$. CF=0, так как $128<256$; OF=1, так как $127+1$ выходит за знаковый диапазон. ZF=0, SF=1.

\subsection*{З20. Заём не равен переносу сложения}\label{sol:z20}
\hyperref[task:z20]{К условию.}
Для $35-60$: код 35 равен $23_{16}$; код $-60$ равен $C4_{16}$ (196). Сумма $35+196=231=E7_{16}$, слово \texttt{11100111}. Unsigned 231, signed $-25$.
Для SUB $35<60$, поэтому CF=1. Точное $-25$ представимо, OF=0; ZF=0, SF=1. Переноса вспомогательного сложения нет, но это не CF вычитания.

Для $(-110)-30$: уменьшаемое имеет код $92_{16}$ (146), дополнительный код $-30$ равен $E2_{16}$ (226). Вспомогательная сумма $146+226=372$, младший байт $74_{16}$ (116 в обеих интерпретациях).
При исходном вычитании сравниваются беззнаковые коды $146$ и $30$: займа нет, CF=0. Точная знаковая разность $-140$ не помещается: OF=1. ZF=0, SF=0.

\section{З21--З24. Операции и маски}
\subsection*{З21. Побитовое сопоставление}\label{sol:z21}
\hyperref[task:z21]{К условию.}
Исходные слова \texttt{10101100} и \texttt{01101001}.
AND оставляет общие единицы: \texttt{00101000}, то есть $28_{16}$.
OR оставляет единицу везде, где она есть хотя бы в одном слове: \texttt{11101101}, то есть $ED_{16}$.
XOR оставляет единицы только там, где биты различны: \texttt{11000101}, то есть $C5_{16}$.
NOT первого слова: \texttt{01010011}, то есть $53_{16}$.
Проверка инверсии: $172+83=255$.

\subsection*{З22. Пять масок}\label{sol:z22}
\hyperref[task:z22]{К условию.}
$53_{16}=\texttt{01010011}$. Все действия независимы:
\begin{enumerate}
    \item Бит 7: OR с $80_{16}$ даёт $D3_{16}=\texttt{11010011}$.
    \item Сброс бита 4: AND с $EF_{16}=\texttt{11101111}$ даёт $43_{16}=\texttt{01000011}$.
    \item Инверсия бита 0: XOR с $01_{16}$ даёт $52_{16}=\texttt{01010010}$.
    \item Проверка бита 6: AND с $40_{16}$ даёт $40_{16}=64$. Результат ненулевой, следовательно, бит установлен.
    \item Младшие 3 бита: AND с $07_{16}$ даёт $03_{16}=\texttt{00000011}$.
\end{enumerate}

\subsection*{З23. Извлечение поля и сборка байта}\label{sol:z23}
\hyperref[task:z23]{К условию.}
$DA_{16}=\texttt{11011010}$. Чтобы исходные биты 4--6 оказались в позициях 0--2, логически сдвигаем вправо на 4: получаем $0D_{16}$. Затем AND с $07_{16}$ убирает исходный бит 7: результат $05_{16}=5$.

Для сброса бита 3 используем маску $F7_{16}=\texttt{11110111}$:
\[
DA_{16}\mathbin{\mathrm{AND}}F7_{16}=D2_{16}.
\]

Сборка другого байта: $A_{16}$ сдвигаем влево на 4, получаем $A0_{16}$. OR с $05_{16}$ даёт $A5_{16}=\texttt{10100101}$.

\subsection*{З24. Сначала очистить, затем установить}\label{sol:z24}
\hyperref[task:z24]{К условию.}
AND с маской $F0_{16}=\texttt{11110000}$ сохраняет старшую тетраду и очищает младшую:
$B6_{16}\mathbin{\mathrm{AND}}F0_{16}=B0_{16}$.
Затем $B0_{16}\mathbin{\mathrm{OR}}09_{16}=B9_{16}$.

OR сохраняет старые единицы младшей тетрады: \texttt{0110} и \texttt{1001} дают \texttt{1111}, поэтому результат равен $BF_{16}$. Для замены поля сначала очистите его маской.

\section{З25--З28. Сдвиги и разрядность}
\subsection*{З25. Потеря старших битов}\label{sol:z25}
\hyperref[task:z25]{К условию.}
$6D_{16}=\texttt{01101101}=109$.
\begin{itemize}
    \item Влево на 1: \texttt{11011010}, $DA_{16}=218$. Точное $109\cdot2=218$ помещается, потерянный старший бит был нулём.
    \item Влево на 3: \texttt{01101000}, $68_{16}=104$. Точное $109\cdot8=872$ не помещается; $872=3\cdot256+104$.
    \item Вправо на 2: \texttt{00011011}, $1B_{16}=27$. Проверка: $\lfloor109/4\rfloor=27$.
\end{itemize}
Левый сдвиг на 1 корректен как unsigned-умножение, хотя старший бит результата 1. Знаковость надо задавать отдельно.

\subsection*{З26. Округление отрицательных чисел}\label{sol:z26}
\hyperref[task:z26]{К условию.}
$93_{16}=147$, знаковое значение $147-256=-109$.
Исходные биты \texttt{10010011}.
\begin{itemize}
    \item Арифметический вправо на 1: \texttt{11001001}, $C9_{16}$, значение $201-256=-55=\lfloor-109/2\rfloor$.
    \item Арифметический вправо на 3: \texttt{11110010}, $F2_{16}$, значение $242-256=-14=\lfloor-109/8\rfloor$.
    \item Логический вправо на 3: \texttt{00010010}, $12_{16}=18=\lfloor147/8\rfloor$.
\end{itemize}
Точное частное $-109/8=-13.625$. При округлении к нулю получаем $-13$, при округлении вниз --- $-14$. Поэтому деление в C и арифметический сдвиг дали разные ответы.

\subsection*{З27. Возврат вытесненных битов}\label{sol:z27}
\hyperref[task:z27]{К условию.}
$81_{16}=\texttt{10000001}$. При циклическом сдвиге влево на 1 старшая единица возвращается в младший бит: \texttt{00000011}, $03_{16}$.

При циклическом сдвиге вправо на 2 отделяем младшие \texttt{01}, переносим их в начало, а оставшиеся шесть битов сдвигаем: \texttt{01 100000}, то есть \texttt{01100000}, $60_{16}$.

При логическом сдвиге влево на 1 старшая единица пропала бы: \texttt{00000010}, $02_{16}$.

\subsection*{З28. Сохранение значения}\label{sol:z28}
\hyperref[task:z28]{К условию.}
$D3_{16}=211$ как unsigned и $211-256=-45$ как дополнительный код.
Нулевое расширение: \texttt{00000000 11010011}, $00D3_{16}$, значение 211.
Знаковое: \texttt{11111111 11010011}, $FFD3_{16}$; беззнаковое значение этого кода 65491, поэтому signed $65491-65536=-45$.

$0137_{16}=256+3\cdot16+7=311$. После сужения остаётся $37_{16}=55$. Значение потеряно: $311>255$, а $311\bmod256=55$.

\section{З29--З32. Объединяем навыки}
\subsection*{З29. Интерпретация, расширение, сложение}\label{sol:z29}
\hyperref[task:z29]{К условию.}
$E9_{16}=233$, знаковое $233-256=-23$. Старший бит равен 1, поэтому до 16 битов расширяем единицами: \texttt{11111111 11101001}, $FFE9_{16}$.

Прибавляем 8 к исходному байту: $E9_{16}+08_{16}=F1_{16}$. Получаем \texttt{11110001}; $241-256=-15$, что совпадает с $-23+8$. CF=0, поскольку $233+8<256$; OF=0, поскольку $-15$ представимо; ZF=0, SF=1.

\subsection*{З30. Три неверных обобщения}\label{sol:z30}
\hyperref[task:z30]{К условию.}
\begin{enumerate}
    \item В 4-битном дополнительном коде \texttt{1111} действительно означает $-1$, но как unsigned это 15, в прямом коде $-7$, в обратном --- отрицательный ноль. Необходимо знать интерпретацию.
    \item Контрпример в 8 битах: $FF_{16}+01_{16}$ даёт $00_{16}$ с CF=1. Знаково это $-1+1=0$, поэтому OF=0. CF не определяет OF.
    \item В 8 битах $0F_{16}=\texttt{00001111}$. Инвертируются все восемь битов, результат \texttt{11110000}, то есть $F0_{16}$. Нули исходной записи тоже участвуют в операции.
\end{enumerate}

\subsection*{З31. Первый неверный шаг}\label{sol:z31}
\hyperref[task:z31]{К условию.}
Код $-3$ указан правильно: $256-3=253=FD_{16}=\texttt{11111101}$.
Ошибка начинается с отождествления арифметического сдвига и деления с округлением к нулю.

Сдвигаем биты вправо, дополняя слева единицей: \texttt{11111110}, код $FE_{16}$, значение $254-256=-2$.
Точное $-3/2=-1.5$ округляется вниз до $-2$, а к нулю --- до $-1$. Если бы исходное число делилось на 2 без остатка, этого расхождения не было бы.

\subsection*{З32. Необычная разрядность}\label{sol:z32}
\hyperref[task:z32]{К условию.}
В 7 битах коды $0\ldots127$. После вычитания смещения 64 диапазон значений $[-64,63]$.

$u_A=-25+64=39=\texttt{0100111}$;
$u_B=38+64=102=\texttt{1100110}$.
Точная сумма $A+B=13$ представима. Правильный смещённый код:
\[
u_{A+B}=39+102-64=77=\texttt{1001101}.
\]
Проверка: $77-64=13$.

Если просто сложить коды, получится 141. После усечения до 7 битов останется $141-128=13$, слово \texttt{0001101}. В исходном смещённом формате это $13-64=-51$, неверная сумма. Но в 7-битном дополнительном коде то же слово \texttt{0001101} правильно представляет положительное число 13. Значение битового слова зависит от формата представления.

\section{Контрольная по занятиям №1--2}\label{solutions:mock}
Оценивайте каждый подпункт по приведённым компонентам. Баллы за верный метод можно получить при отдельной арифметической ошибке, если дальнейшие действия последовательно выполнены над своим результатом. За компонент начисляется не больше указанного максимума. Общая сумма --- 100 баллов.

\subsection*{К1. Переводы}\label{sol:k1}
\hyperref[task:k1]{К условию.}
$173=128+32+8+4+1$, поэтому $10101101_2$. Группы по три: \texttt{010 101 101}, ответ $255_8$; по четыре: \texttt{1010 1101}, ответ $AD_{16}$. Проверки: $2\cdot64+5\cdot8+5=173$ и $10\cdot16+13=173$.

$2E6_{16}$ заменяем тетрадами \texttt{0010 1110 0110}. Получаем $1011100110_2$; триады \texttt{001 011 100 110} дают $1346_8$. Проверка: $2\cdot256+14\cdot16+6=742=512+3\cdot64+4\cdot8+6$.

\textbf{Баллы:} подпункт 1 --- 2 за способ перевода, по 1 за каждую из трёх записей; подпункт 2 --- 2 за правильные тетрады, 1 за двоичную запись, 1 за перегруппировку, 1 за восьмеричный ответ.

\subsection*{К2. Арифметика}\label{sol:k2}
\hyperref[task:k2]{К условию.}
Складываем $3A6_{16}+07B_{16}$. Младший разряд: $6+11=17$, цифра 1, перенос 1. Средний: $10+7+1=18$, цифра 2, перенос 1. Старший: $3+0+1=4$. Ответ $421_{16}$; проверка $934+123=1057$.

При делении $264_8:7_8$ берём $26_8$. Произведение $7_8\cdot3_8=25_8$, остаток 1. Сносим 4: $14_8$. Следующая цифра частного 1, остаток $14_8-7_8=5_8$. Ответ $31_8$, остаток $5_8$. Проверка: $180=7\cdot25+5$ и $0\le5<7$.

\textbf{Баллы:} сложение --- 2 за разрядные действия с переносами, 2 за ответ, 1 за проверку; деление --- 2 за подбор цифр, по 1 за частное и остаток, 1 за проверку.

\subsection*{К3. Диапазоны}\label{sol:k3}
\hyperref[task:k3]{К условию.}
10-битный дополнительный код: $[-2^9,2^9-1]=[-512,511]$, количество $2^{10}=1024$.

Для unsigned до 1500: в 10 битах максимум 1023, в 11 --- 2047. Значит 11 битов достаточно и необходимо.

Для $[-600,600]$: 10-битный диапазон $[-512,511]$ недостаточен; 11-битный $[-1024,1023]$ подходит. Ответ 11 битов.

\textbf{Баллы:} подпункт 1 --- по 1 за каждую границу, 2 за число значений; подпункты 2 и 3 --- по 1 за ответ, достаточность и недостаточность меньшей ширины.

\subsection*{К4. Представления}\label{sol:k4}
\hyperref[task:k4]{К условию.}
$53=32+16+4+1=\texttt{00110101}$.
\begin{itemize}
    \item Прямой: \texttt{10110101}, $B5_{16}$.
    \item Обратный: \texttt{11001010}, $CA_{16}$.
    \item Дополнительный: \texttt{11001011}, $CB_{16}$; проверка $256-53=203$.
    \item Со смещением: $-53+128=75$, \texttt{01001011}, $4B_{16}$.
\end{itemize}
$D9_{16}=13\cdot16+9=217$ как unsigned; как дополнительный код $217-256=-39$. Знаковое расширение повторяет единицу: \texttt{11111111 11011001}, $FFD9_{16}$.

\textbf{Баллы:} за каждый из четырёх кодов --- 1 за правильный способ, 1 за согласованные двоичную и hex-записи; за каждую интерпретацию $D9_{16}$ --- 1 за расчёт, 1 за значение; расширение --- 1 за правило и 2 за правильные записи результата.

\subsection*{К5. Флаги}\label{sol:k5}
\hyperref[task:k5]{К условию.}
\begin{enumerate}
    \item $3F_{16}+41_{16}=80_{16}$. Сохранённое знаковое значение $-128$; точная сумма 128. CF=0 ($128<256$), OF=1 ($128>127$), ZF=0, SF=1.
    \item Коды $-90$ и $-50$ равны $A6_{16}$ (166) и $CE_{16}$ (206). Сумма кодов 372, сохраняем $74_{16}$, значение 116. CF=1 ($372\ge256$), OF=1 (точное $-140<-128$), ZF=0, SF=0.
    \item $25-70=-45$, код $256-45=211=D3_{16}$. CF=1 ($25<70$), OF=0 (разность представима), ZF=0, SF=1.
    \item Код $-100$ равен $9C_{16}$ (156). Вычитание кода 40 даёт $156-40=116=74_{16}$, знаковое значение 116. CF=0 ($156\ge40$), OF=1 (точное $-140$ непредставимо), ZF=0, SF=0.
\end{enumerate}
В подпунктах 2 и 4 точный знаковый ответ и сохранённые биты совпадают, но CF различается: выполнялись разные операции над разными кодами операндов.

\textbf{Баллы за каждый подпункт:} 1 за вычисление кода, 1 за знаковое значение, 1 за CF с обоснованием, 1 за OF с обоснованием, по 0.5 за ZF и SF. Всего $4\cdot5=20$.

\subsection*{К6. Маски}\label{sol:k6}
\hyperref[task:k6]{К условию.}
$CA_{16}=\texttt{11001010}$, $75_{16}=\texttt{01110101}$.
\begin{itemize}
    \item AND: \texttt{01000000}, $40_{16}$.
    \item OR: \texttt{11111111}, $FF_{16}$.
    \item XOR: \texttt{10111111}, $BF_{16}$.
    \item NOT $x$: \texttt{00110101}, $35_{16}$.
\end{itemize}
В бите 6 у обоих слов единица: AND сохраняет её, XOR обнуляет. В остальных позициях единицы операндов не совпадают.

Замена тетрады:
\[
(CA_{16}\mathbin{\mathrm{AND}}F0_{16})\mathbin{\mathrm{OR}}03_{16}
=C0_{16}\mathbin{\mathrm{OR}}03_{16}=C3_{16}.
\]
\textbf{Баллы:} за каждую из четырёх операций --- 1 за правильные биты, 1 за hex; замена поля --- 1 за маску и операции, 1 за результат. Всего 10.

\subsection*{К7. Сдвиги и расширение}\label{sol:k7}
\hyperref[task:k7]{К условию.}
Исходные биты $A9_{16}=\texttt{10101001}$, unsigned 169, signed $169-256=-87$.
\begin{enumerate}
    \item Логический вправо на 2: \texttt{00101010}, $2A_{16}=42$. Проверка $\lfloor169/4\rfloor=42$.
    \item Арифметический вправо на 2: \texttt{11101010}, $EA_{16}$; $234-256=-22$. Проверка $\lfloor-87/4\rfloor=-22$.
    \item Циклический влево на 1: вытесненная единица возвращается справа. Получаем \texttt{01010011}, $53_{16}$.
    \item Знаковое расширение: \texttt{11111111 10101001}, $FFA9_{16}$.
    \item Деление с округлением к нулю: $-87/4=-21.75\longrightarrow-21$. Это на единицу больше результата арифметического сдвига, округляющего вниз.
\end{enumerate}
\textbf{Баллы:} в подпунктах 1 и 2 по 1 за биты, hex и значение; в подпунктах 3 и 4 по 1 за правило, биты и hex; в подпункте 5 --- по 1 за исходное знаковое значение, частное и объяснение различия. Всего 15.

\subsection*{К8. Исправление решения}\label{sol:k8}
\hyperref[task:k8]{К условию.}
Полная сумма кодов действительно $FD_{16}+03_{16}=100_{16}$, то есть 256. Но в 8-битном результате остаётся $00_{16}$. Знаково исходные операнды $253-256=-3$ и 3, точная сумма 0.

\begin{itemize}
    \item CF=1: есть перенос за 8 битов. Эта часть исходного решения верна.
    \item OF=0: точная знаковая сумма 0 представима. Вывод OF из CF неверен.
    \item ZF=1: сохранённый байт равен нулю.
    \item SF=0: старший бит сохранённого результата нулевой. Знак первого операнда сам по себе не определяет SF.
\end{itemize}
\textbf{Баллы:} 2 за правильное усечение до байта; 2 за интерпретацию слагаемых и знаковую сумму; по 1 за каждый флаг (всего 4); 2 за объяснение ошибок в определении OF, ZF и SF. Итого 10.

\section{Повторение по занятиям №1--2}
Материал для повторения приведён в основном пособии \href{adaptation_course.pdf}{«Адаптационный курс»}. В таблице указаны соответствующие занятия и задания практикума.

\begin{center}
\small
\begin{tabular}{>{\raggedright\arraybackslash}p{0.31\textwidth}|>{\raggedright\arraybackslash}p{0.35\textwidth}|>{\raggedright\arraybackslash}p{0.23\textwidth}}
Ошибка & Материал курса & Задания практикума\\\hline
Неверный порядок цифр, потерянные нули & \href{adaptation_course.pdf}{Лекция и семинар №1}: переводы между системами счисления & \hyperref[ex:p1]{П1}, \hyperref[ex:p2]{П2}; З1--З4\\[4pt]
Неверный перенос, заём или остаток & \href{adaptation_course.pdf}{Лекция и семинар №1}: целочисленная арифметика & \hyperref[ex:p3]{П3}; З5--З8\\[4pt]
Ошибка на границе диапазона & \href{adaptation_course.pdf}{Лекция и семинар №2}: диапазоны целых чисел & \hyperref[ex:p4]{П4}; З9--З12\\[4pt]
Код принят за модуль или значение & \href{adaptation_course.pdf}{Лекция и семинар №2}: коды целых чисел & \hyperref[ex:p5]{П5}, \hyperref[ex:p6]{П6}; З13--З16\\[4pt]
CF и OF смешиваются & \href{adaptation_course.pdf}{Лекция и семинар №2}: арифметика в дополнительном коде, флаги & \hyperref[ex:p7]{П7}, \hyperref[ex:p8]{П8}; З17--З20\\[4pt]
Маска меняет лишние биты & \href{adaptation_course.pdf}{Лекция и семинар №2}: побитовые операции & \hyperref[ex:p9]{П9}; З21--З24\\[4pt]
Сдвиг даёт неверный знак или частное & \href{adaptation_course.pdf}{Лекция и семинар №2}: сдвиги и знаковое расширение & \hyperref[ex:p10]{П10}; З25--З28\\[4pt]
Ошибки в цепочке вычислений & \href{adaptation_course.pdf}{Лекции и семинары №1--2}: переводы, представление чисел и арифметика & З29--З32, затем К8
\end{tabular}
\end{center}
